% AA-59
\chapter[Podalgebry Cartana]{Podalgebry Cartana i rozkład na przestrzenie korzeniowe}
\label{aa:ch:cartan-root-decomposition}

W rozdziale~\ref{aa:ch:semisimple-lie} wykazaliśmy niezdegenerowanie formy Killinga algebry półprostej oraz wewnętrzność jej derywacji. Teraz wyprowadzimy z tych wyników rozkład algebry na wspólne przestrzenie własne operatorów reprezentacji dołączonej. W całym rozdziale $K$ jest algebraicznie domkniętym ciałem charakterystyki zero, a $\mathfrak g$ jest skończenie wymiarową algebrą Liego nad $K$.

\section{Rozkład Jordana derywacji}

\begin{lemat}[Uogólnione przestrzenie własne derywacji]
\label{aa:lem:derivation-generalized-weights}
Niech $\mathfrak g$ będzie skończenie wymiarową algebrą Liego nad algebraicznie domkniętym ciałem $K$ charakterystyki zero, a $D$ jej derywacją. Dla każdego $\lambda\in K$ oznaczmy przez $V_\lambda$ uogólnioną przestrzeń własną $D$, czyli $\bigcup_{m\ge1}\ker(D-\lambda I)^m$ (równą zeru, jeśli $\lambda$ nie jest wartością własną). Wówczas dla wszystkich $\lambda,\mu\in K$ zachodzi $[V_\lambda,V_\mu]\subset V_{\lambda+\mu}$.
\end{lemat}
\begin{proof}
Dla $x\in V_\lambda$, $y\in V_\mu$ z tożsamości derywacji otrzymujemy
\[
 (D-(\lambda+\mu)I)[x,y]=
 [(D-\lambda I)x,y]+[x,(D-\mu I)y].
\]
Indukcja względem $m$ daje
\[
 (D-(\lambda+\mu)I)^m[x,y]=
 \sum_{j=0}^m\binom mj[(D-\lambda I)^jx,(D-\mu I)^{m-j}y].
\]
W kroku indukcyjnym stosujemy poprzednią tożsamość do każdego nawiasu; po zmianie indeksu współczynniki sumują się według wzoru Pascala. Wybieramy $p,q\ge1$ takie, że $(D-\lambda I)^px=0$ i $(D-\mu I)^qy=0$. Dla $m=p+q-1$ każdy składnik znika: albo $j\ge p$, albo $m-j\ge q$. Stąd teza, także gdy operator $D-(\lambda+\mu)I$ jest odwracalny.
\end{proof}

\begin{twierdzenie}[Rozkład Jordana w algebrze półprostej]
\label{aa:tw:inner-jordan}
Niech $\mathfrak g$ będzie skończenie wymiarową półprostą algebrą Liego nad algebraicznie domkniętym ciałem $K$ charakterystyki zero. Dla każdego $x\in\mathfrak g$ istnieją jednoznaczne $s,n\in\mathfrak g$ takie, że $x=s+n$, $[s,n]=0$, $\operatorname{ad}_s$ jest diagonalizowalny, a $\operatorname{ad}_n$ nilpotentny. Dla każdego $y\in\mathfrak g$, jeśli $[x,y]=0$, to $[s,y]=[n,y]=0$.
\end{twierdzenie}
\begin{proof}
Niech $D=\operatorname{ad}_x$. Rozkład na uogólnione przestrzenie własne daje $\mathfrak g=\bigoplus_\lambda V_\lambda$. Definiujemy $S|_{V_\lambda}=\lambda I$ oraz $N=D-S$. Operator $S$ jest diagonalizowalny, $N$ nilpotentny, a oba komutują, gdyż zachowują wszystkie $V_\lambda$ i $S$ jest na nich skalarny. Na mocy Lematu~\ref{aa:lem:derivation-generalized-weights}, dla $u\in V_\lambda$, $v\in V_\mu$ mamy
$S[u,v]=(\lambda+\mu)[u,v]=[Su,v]+[u,Sv]$.
Dwuliniowość i rozkład na sumę prostą rozszerzają tę równość na dowolne $u,v$. Zatem $S$ i $N=D-S$ są derywacjami.

Twierdzenie~\ref{aa:tw:semisimple-derivations} daje $s,n$ z $S=\operatorname{ad}_s$, $N=\operatorname{ad}_n$. Centrum $\mathfrak g$ jest zerowe, więc z $D=S+N$ i $[S,N]=0$ wynikają $x=s+n$ i $[s,n]=0$.

Dla jednoznaczności niech $D=S'+N'$ będzie innym takim rozkładem. Na przestrzeni własnej $W_\nu$ operatora $S'$ operator $N'$ zachowuje $W_\nu$ dzięki komutowaniu, a $D-\nu I=N'$ jest tam nilpotentny. Stąd $W_\nu\subset V_\nu$. Dla $v\in V_\lambda$ zapisujemy $v=\sum_\nu w_\nu$, gdzie $w_\nu\in W_\nu\subset V_\nu$ (przyjmujemy $W_\nu=0$, jeśli $\nu$ nie jest wartością własną $S'$). Jednoznaczność składowych w rozkładzie na $V_\nu$ daje $w_\nu=0$ dla $\nu\ne\lambda$ oraz $v=w_\lambda$. Zatem $V_\lambda\subset W_\lambda$, co wraz z wcześniejszą inkluzją daje równość. Zatem $S'=S$, $N'=N$, a zerowe centrum daje jednoznaczność $s,n$.

Jeśli $[x,y]=0$, operator $\operatorname{ad}_y$ komutuje z $D$, a zatem zachowuje każde $V_\lambda$ (komutuje z każdą potęgą $D-\lambda I$). Komutuje więc z $S$ i $N=D-S$. Ponownie zerowe centrum daje żądane równości.
\end{proof}

\section{Podalgebry toralne i Cartana}

\begin{definicja}
\label{aa:def:toral-cartan}
Niech $\mathfrak g$ będzie skończenie wymiarową algebrą Liego nad algebraicznie domkniętym ciałem $K$ charakterystyki zero. Podalgebra $\mathfrak t\subset\mathfrak g$ jest toralna, jeśli jest abelowa i dla każdego $t\in\mathfrak t$ operator $\operatorname{ad}_t$ na $\mathfrak g$ jest diagonalizowalny. Podalgebra $\mathfrak h$ jest podalgebrą Cartana, jeśli jest nilpotentna i jej normalizator
$N_{\mathfrak g}(\mathfrak h)=\{x\in\mathfrak g:[x,\mathfrak h]\subset\mathfrak h\}$ jest równy $\mathfrak h$. Centralizatorem podalgebry $\mathfrak h$ nazywamy $Z_{\mathfrak g}(\mathfrak h)=\{x\in\mathfrak g:[x,t]=0\ \forall t\in\mathfrak h\}$.
\end{definicja}

\begin{lemat}[Wspólny rozkład własny]
\label{aa:lem:toral-weight}
Niech $\mathfrak g$ będzie skończenie wymiarową algebrą Liego nad algebraicznie domkniętym ciałem $K$ charakterystyki zero, a $\mathfrak t$ jej podalgebrą toralną. Istnieje skończony zbiór $\Lambda\subset\mathfrak t^*$ takich funkcjonałów, że
\[
 \mathfrak g=\bigoplus_{\lambda\in\Lambda}\mathfrak g_\lambda,\qquad
 \mathfrak g_\lambda=\{x:[t,x]=\lambda(t)x\ \forall t\in\mathfrak t\},
\]
gdzie $\Lambda$ zawiera dokładnie te $\lambda$, dla których $\mathfrak g_\lambda\ne0$. Dla wszystkich $\lambda,\mu\in\mathfrak t^*$ (z zerowymi przestrzeniami dopuszczonymi) zachodzą
$[\mathfrak g_\lambda,\mathfrak g_\mu]\subset\mathfrak g_{\lambda+\mu}$ oraz
$\kappa(\mathfrak g_\lambda,\mathfrak g_\mu)=0$, jeśli $\lambda+\mu\ne0$, gdzie $\kappa$ jest formą Killinga.
\end{lemat}
\begin{proof}
Wybieramy bazę $t_1,\ldots,t_r$ przestrzeni $\mathfrak t$. Operatory $\operatorname{ad}_{t_i}$ komutują. Rozkład na przestrzenie własne pierwszego jest zachowany przez każdy pozostały operator. Ograniczenie operatora diagonalizowalnego do niezmienniczej podprzestrzeni pozostaje diagonalizowalne: jego wielomian minimalny dzieli wielomian o parami różnych pierwiastkach. Powtarzając rozkład na każdej składowej, otrzymujemy wspólną bazę własną. Wspólna lista wartości własnych na bazie $\mathfrak t$ wyznacza funkcjonał liniowy $\lambda$; liniowość reprezentacji dołączonej daje równość dla wszystkich $t$. Gdy $\mathfrak t=0$, jedynym możliwym funkcjonałem jest zero.

Ustalmy $\lambda,\mu\in\mathfrak t^*$, $x\in\mathfrak g_\lambda$, $y\in\mathfrak g_\mu$ oraz $t\in\mathfrak t$. Z tożsamości Jacobiego mamy
$[t,[x,y]]=[[t,x],y]+[x,[t,y]]=(\lambda+\mu)(t)[x,y]$.
Z niezmienniczości i symetrii formy Killinga wynika
$0=\kappa([t,x],y)+\kappa(x,[t,y])=(\lambda+\mu)(t)\kappa(x,y)$.
Jeżeli $\lambda+\mu\ne0$, wybieramy $t$, dla którego ten funkcjonał nie znika.
\end{proof}

\begin{twierdzenie}[Maksymalna podalgebra toralna]
\label{aa:tw:maximal-toral-cartan}
Niech $\mathfrak g$ będzie skończenie wymiarową półprostą algebrą Liego nad algebraicznie domkniętym ciałem $K$ charakterystyki zero. Istnieje maksymalna względem inkluzji podalgebra toralna $\mathfrak h$. Dla każdej takiej podalgebry zachodzą
$Z_{\mathfrak g}(\mathfrak h)=N_{\mathfrak g}(\mathfrak h)=\mathfrak h$, a ograniczenie formy Killinga do $\mathfrak h$ jest niezdegenerowane. W szczególności $\mathfrak h$ jest podalgebrą Cartana.
\end{twierdzenie}
\begin{proof}
Podalgebra zerowa jest toralna. Wybierając toralną podalgebrę największego wymiaru, otrzymujemy podalgebrę maksymalną. Niech teraz $\mathfrak h$ będzie dowolną maksymalną toralną podalgebrą. Jej centralizator $\mathfrak c=\mathfrak g_0$ jest podalgebrą (tożsamość Jacobiego). Na mocy kryterium półprostoty Cartana (Twierdzenie~\ref{aa:tw:cartan-semisimple}) forma Killinga $\kappa=\kappa_{\mathfrak g}$ jest niezdegenerowana. Wykażemy, że jej ograniczenie do $\mathfrak c$ również jest niezdegenerowane. Jeśli $u\in\mathfrak c$ i $\kappa(u,v)=0$ dla każdego $v\in\mathfrak c$, to na mocy Lematu~\ref{aa:lem:toral-weight} mamy także $\kappa(u,v)=0$ dla $v\in\mathfrak g_\lambda$, $\lambda\ne0$. Każdy element $\mathfrak g$ jest sumą składowej z $\mathfrak c$ i składowych z niezerowych przestrzeni wagowych; dwuliniowość daje więc $\kappa(u,\mathfrak g)=0$. Niezdegenerowanie $\kappa$ oznacza $u=0$. Chodzi tu o ograniczenie formy Killinga całej $\mathfrak g$, a nie o formę Killinga samej $\mathfrak c$.

Dla $x\in\mathfrak c$ rozkładamy $x=s+n$ według Twierdzenia~\ref{aa:tw:inner-jordan}. Obie składowe komutują z $\mathfrak h$. Przestrzeń $\mathfrak h+Ks$ jest abelową podalgebrą; każdy jej element działa diagonalizowalnie, gdyż komutujące operatory diagonalizowalne $\operatorname{ad}_s$ i $\operatorname{ad}_t$ dla $t\in\mathfrak h$ mają wspólną bazę własną. Maksymalność daje $s\in\mathfrak h$. Wobec tego $\operatorname{ad}_x|_{\mathfrak c}=\operatorname{ad}_n|_{\mathfrak c}$ jest nilpotentny. Twierdzenie Engela (Twierdzenie~\ref{aa:tw:engel}) zastosowane do $\mathfrak c$ daje jej nilpotentność, a więc rozwiązalność na mocy Lematu~\ref{aa:lem:lie-series}.

Na mocy Wniosku~\ref{aa:cor:lie-flag} reprezentacja $\mathfrak c\to\mathfrak{gl}(\mathfrak g)$, $u\mapsto\operatorname{ad}_u$ (ograniczenie reprezentacji dołączonej $\mathfrak g$) ma wspólną bazę górnotrójkątną. Operator $\operatorname{ad}_n$ jest nilpotentny, więc jego przekątna w tej bazie jest zerowa. Dla każdego $y\in\mathfrak c$ iloczyn $\operatorname{ad}_n\operatorname{ad}_y$ jest ściśle górnotrójkątny, zatem $\kappa(n,y)=0$. Niezdegenerowanie ograniczenia do $\mathfrak c$ i $n\in\mathfrak c$ dają $n=0$. Stąd $x=s\in\mathfrak h$ i $\mathfrak c=\mathfrak h$.

Niech $x\in N_{\mathfrak g}(\mathfrak h)$. Zapiszmy $x=x_0+\sum_{\lambda\ne0}x_\lambda$ według Lematu~\ref{aa:lem:toral-weight}. Dla każdego $t\in\mathfrak h$ mamy $[t,x]=\sum_{\lambda\ne0}\lambda(t)x_\lambda\in\mathfrak h=\mathfrak g_0$. Jednoznaczność składowych i wybór $t$ z $\lambda(t)\ne0$ dają $x_\lambda=0$ dla każdej $\lambda\ne0$. Zatem $x\in\mathfrak h$. Odwrotna inkluzja wynika z abelowości, która daje też nilpotentność $\mathfrak h$.
\end{proof}

\begin{uwaga}
W tym rozdziale pracujemy z podalgebrami Cartana otrzymanymi jako maksymalne podalgebry toralne. Odwrotność Twierdzenia~\ref{aa:tw:maximal-toral-cartan} oraz sprzężoność podalgebr Cartana nie są tu potrzebne i pozostają do osobnego dowodu.
\end{uwaga}

\section{Przestrzenie korzeniowe i forma Killinga}

\begin{definicja}
\label{aa:def:lie-root}
Niech $\mathfrak g$ będzie skończenie wymiarową półprostą algebrą Liego nad algebraicznie domkniętym ciałem $K$ charakterystyki zero, a $\mathfrak h$ jej maksymalną podalgebrą toralną. Korzeniem nazywamy niezerowy funkcjonał $\alpha\in\mathfrak h^*$ taki, że $\mathfrak g_\alpha\ne0$, gdzie $\mathfrak g_\alpha$ określono w Lemacie~\ref{aa:lem:toral-weight}. Zbiór korzeni oznaczamy $\Phi$. Rozkład korzeniowy ma postać
\[
 \mathfrak g=\mathfrak h\oplus\bigoplus_{\alpha\in\Phi}\mathfrak g_\alpha.
\]
Wynika on z Twierdzenia~\ref{aa:tw:maximal-toral-cartan}.
\end{definicja}

\begin{twierdzenie}
\label{aa:tw:root-killing-pairing}
Niech $\mathfrak g$ będzie skończenie wymiarową półprostą algebrą Liego nad algebraicznie domkniętym ciałem $K$ charakterystyki zero, $\mathfrak h$ jej maksymalną podalgebrą toralną, a $\Phi$ zbiorem korzeni względem $\mathfrak h$ z Definicji~\ref{aa:def:lie-root}. Korzenie rozpinają $\mathfrak h^*$. Dla każdego $\alpha\in\Phi$ mamy $-\alpha\in\Phi$, a forma Killinga $\kappa$ algebry $\mathfrak g$ określa niezdegenerowane parowanie $\mathfrak g_\alpha\times\mathfrak g_{-\alpha}\to K$. W szczególności $\dim\mathfrak g_\alpha=\dim\mathfrak g_{-\alpha}$.
\end{twierdzenie}
\begin{proof}
Jeśli $t\in\mathfrak h$ znika pod każdym korzeniem, to $\operatorname{ad}_t$ znika na wszystkich przestrzeniach korzeniowych oraz na $\mathfrak h$. Jest więc zerowy i $t=0$ dzięki zerowemu centrum. Gdyby rozpiętość korzeni była właściwą podprzestrzenią $\mathfrak h^*$, istniałby niezerowy $t$ w jej anihilatorze (uzupełniamy bazę tej podprzestrzeni do bazy $\mathfrak h^*$). Sprzeczność.

Dla $0\ne x\in\mathfrak g_\alpha$ niezdegenerowanie $\kappa$ daje $y\in\mathfrak g$ z $\kappa(x,y)\ne0$. Rozkład $y$ i Lemat~\ref{aa:lem:toral-weight} pokazują, że jedyny składnik mogący dać niezerową wartość należy do $\mathfrak g_{-\alpha}$. Dowodzi to istnienia przeciwnego korzenia i zerowości lewego jądra parowania. Stosując argument do $-\alpha$, otrzymujemy zerowość prawego jądra. Iniekcje do odpowiednich przestrzeni dualnych dają obie nierówności wymiarów.
\end{proof}

\begin{twierdzenie}[Trójka typu $\mathfrak{sl}_2$]
\label{aa:tw:root-sl2-triple}
Niech $\mathfrak g$ będzie skończenie wymiarową półprostą algebrą Liego nad algebraicznie domkniętym ciałem $K$ charakterystyki zero, $\mathfrak h$ jej maksymalną podalgebrą toralną, a $\Phi$ zbiorem korzeni względem $\mathfrak h$ z Definicji~\ref{aa:def:lie-root}. Niech $\kappa$ oznacza formę Killinga $\mathfrak g$. Dla $\alpha\in\Phi$ niech $t_\alpha\in\mathfrak h$ będzie jedynym elementem takim, że $\kappa(t_\alpha,t)=\alpha(t)$ dla wszystkich $t\in\mathfrak h$. Wtedy $\alpha(t_\alpha)\ne0$. Dla każdego $0\ne e\in\mathfrak g_\alpha$ istnieje $f\in\mathfrak g_{-\alpha}$ takie, że dla $h_\alpha=2t_\alpha/\alpha(t_\alpha)$ zachodzą
\[
 [h_\alpha,e]=2e,\quad [h_\alpha,f]=-2f,\quad [e,f]=h_\alpha.
\]
Elementy $e,f,h_\alpha$ rozpinają podalgebrę izomorficzną z $\mathfrak{sl}_2(K)$.
\end{twierdzenie}
\begin{proof}
Istnienie i jednoznaczność $t_\alpha$ wynikają z niezdegenerowania $\kappa|_{\mathfrak h\times\mathfrak h}$. Dla $x\in\mathfrak g_\alpha$, $y\in\mathfrak g_{-\alpha}$ mamy $[x,y]\in\mathfrak g_0=\mathfrak h$. Dla $t\in\mathfrak h$ niezmienniczość daje
$\kappa([x,y],t)=\kappa(x,[y,t])=\alpha(t)\kappa(x,y)$.
Niezdegenerowanie na $\mathfrak h$ daje tożsamość $[x,y]=\kappa(x,y)t_\alpha$.

Ustalmy $0\ne x\in\mathfrak g_\alpha$ i wybierzmy $y\in\mathfrak g_{-\alpha}$ z $\kappa(x,y)=1$, co jest możliwe dzięki Twierdzeniu~\ref{aa:tw:root-killing-pairing} i przeskalowaniu. Wtedy $[x,y]=t_\alpha$. Załóżmy, że $\alpha(t_\alpha)=0$. Wówczas $t_\alpha$ komutuje z $x$ i $y$, więc operatory $\operatorname{ad}_x,\operatorname{ad}_y$ zachowują wszystkie przestrzenie własne diagonalizowalnego operatora $T=\operatorname{ad}_{t_\alpha}$. Na każdej niezerowej przestrzeni własnej $W$ o wartości własnej $\nu$ zachodzi
\[
 [\operatorname{ad}_x|_W,\operatorname{ad}_y|_W]=\nu I_W.
\]
Ślad komutatora jest zerowy, więc $\nu\dim W=0$, a charakterystyka zero daje $\nu=0$. Z diagonalizowalności wynika $T=0$ i z zerowego centrum $t_\alpha=0$, co przeczy $\alpha\ne0$.

Dla ustalonego $e$ wybieramy $y$ z $\kappa(e,y)=1$ i kładziemy $f=2y/\alpha(t_\alpha)$. Tożsamość na nawias daje $[e,f]=h_\alpha$, a definicje przestrzeni korzeniowych dają dwie pozostałe równości. Trzy elementy są niezerowe i należą do różnych składowych sumy prostej (charakterystyka zero daje $\alpha\ne-\alpha$), więc są liniowo niezależne. Równości w tezie wraz z antysymetrią i $[v,v]=0$ dla każdego $v$ pokazują, że ich rozpiętość jest zamknięta na nawias. Ich nawiasy są dokładnie nawiasami standardowej bazy $\mathfrak{sl}_2(K)$ z Przykładu~\ref{aa:prz:heisenberg-sl2}. Odwzorowanie tej bazy na $h_\alpha,e,f$ jest zatem izomorfizmem.
\end{proof}

\section{Przykład macierzowy}

\begin{przyklad}[$\mathfrak{sl}_3(K)$]
\label{aa:prz:sl3-root-decomposition}
Niech $\mathfrak g=\mathfrak{sl}_3(K)$, a $\mathfrak h$ będzie przestrzenią macierzy diagonalnych o śladzie zero. Algebra $\mathfrak g$ jest półprosta: z Twierdzenia~\ref{aa:tw:classical-killing} forma Killinga wynosi $6\operatorname{tr}(XY)$. Jeżeli $X$ jest ortogonalna do wszystkich $E_{ij}$ dla $i\ne j$, to wszystkie jej wpisy pozadiagonalne znikają. Ortogonalność do $E_{ii}-E_{jj}$ wymusza równość wpisów diagonalnych; zerowy ślad i charakterystyka zero dają $X=0$. Stosujemy Twierdzenie~\ref{aa:tw:cartan-semisimple}.

Dla $H=\operatorname{diag}(a_1,a_2,a_3)$, gdzie $\sum_i a_i=0$, mamy
$[H,E_{ij}]=(a_i-a_j)E_{ij}$. Macierze diagonalne komutują, a sześć jednostek pozadiagonalnych wraz z bazą $\mathfrak h$ tworzy wspólną bazę własną. Stąd $\mathfrak h$ jest toralna. Jeżeli macierz $X$ komutuje ze wszystkimi $H$, dla każdego $i\ne j$ wybieramy $H$ z $a_i\ne a_j$; równość $[H,X]_{ij}=0$ daje $X_{ij}=0$. Centralizator jest zatem równy $\mathfrak h$, co wyklucza większą abelową podalgebrę.

Korzeniami są $\alpha_{ij}(H)=a_i-a_j$ dla $i\ne j$, a $\mathfrak g_{\alpha_{ij}}=KE_{ij}$. Funkcjonały są parami różne: ich wektory współczynników mają sumę zero i dwa takie wektory określają ten sam funkcjonał na płaszczyźnie $\sum_i a_i=0$ tylko wtedy, gdy ich różnica jest wielokrotnością $(1,1,1)$, a więc jest zerowa. Dla dowolnego $H\in\mathfrak h$ mamy $\kappa((E_{ii}-E_{jj})/6,H)=\operatorname{tr}((E_{ii}-E_{jj})H)=a_i-a_j$. Z jednoznaczności elementu reprezentującego funkcjonał względem formy Killinga dostajemy
$t_{\alpha_{ij}}=(E_{ii}-E_{jj})/6$ i $\alpha_{ij}(t_{\alpha_{ij}})=1/3$, zatem $h_{\alpha_{ij}}=E_{ii}-E_{jj}$. Standardowa trójka to $E_{ij},E_{ji},h_{\alpha_{ij}}$.

W modelu rzeczywistym płaszczyzny współczynników o sumie zero wektory odpowiadające korzeniom tworzą sześciokąt foremny. Rysunek przedstawia układ typu $A_2$ z rozdziału~\ref{aa:ch:root-systems}; nie utożsamia dowolnego ciała $K$ z podciałem $\mathbb R$.
\begin{center}
\begin{tikzpicture}[scale=1.45, every node/.style={font=\small}]
\foreach \ang/\lab in {0/{\alpha_{12}},60/{\alpha_{13}},120/{\alpha_{23}},180/{\alpha_{21}},240/{\alpha_{31}},300/{\alpha_{32}}}{
 \draw[-{Stealth}] (0,0)--(\ang:1);
 \node at (\ang:1.32) {$\lab$};
}
\fill (0,0) circle (0.025);
\end{tikzpicture}
\end{center}
Na przykład $[E_{12},E_{23}]=E_{13}$, gdyż $E_{12}E_{23}=E_{13}$ i $E_{23}E_{12}=0$. Odpowiada to równości $\alpha_{12}+\alpha_{23}=\alpha_{13}$. Z kolei $[E_{12},E_{13}]=0$, ponieważ obydwa iloczyny są zerowe.
\end{przyklad}

Wykazaliśmy istnienie rozkładu korzeniowego, parowanie przeciwległych przestrzeni oraz istnienie trójek typu $\mathfrak{sl}_2$. Jednowymiarowość przestrzeni korzeniowych, całkowitość liczb Cartana, ciągi korzeniowe i realizacja aksjomatów układu korzeni z rozdziału~\ref{aa:ch:root-systems} wymagają dalszych argumentów z teorii reprezentacji $\mathfrak{sl}_2$; nie zostały tutaj założone ani dowiedzione.
